\chapter{The Cow Came First}

The cow came first.

Father Handy saw it through the sacristy doorway, black-spotted and patient, hauling a metal cart across the dirt yard as if time itself were heavier than the load. Behind it—on the cart—sat Tibor McMasters. Or what remained of him.

A torso. A head. A face that still smiled.

Father Handy stepped back into the church before Tibor noticed him. He needed a moment. Today required words that could not be taken back.

Tibor never saw mornings clearly. They hit him like an illness. Smells, light, motion—everything triggered nausea. It wasn’t weakness; it was refusal. Some part of him resisted waking up in a body that was no longer a body at all.

Father Handy, by contrast, loved mornings. He loved the sun over the Utah pastures, the smell of hot clover, the faint clinking of cow tags in the distance. Life persisted here, stubborn and altered. What troubled him wasn’t the sight of Tibor, but the pain Tibor carried so quietly, like a debt that could never be paid.

Behind the altar waited the mural.

Only a fragment was finished—years of work still ahead. It was meant to last generations. Not eternity, Father Handy corrected himself. Eternity belonged to God. This was human work. Human things decayed.

The church had already adapted. Many worshippers no longer had legs or arms. Genuflection had been officially removed from doctrine.

The cow stopped. Tibor disengaged it and powered the cart forward on battery alone, up the wooden ramp and into the church. Father Handy felt the vibration through the floor, felt Tibor’s arrival like a pressure change in the air.

Tibor was retching, quietly. Fighting it. Always cheerful anyway.

“Do we have coffee?” Father Handy asked his wife.

“Yes,” Ely said. Her voice was dry, dutiful, exhausted. She moved as though devotion were a chore that had worn her down to bone.

“Morning!” Tibor called, bright as ever.

“Black. Hot,” Father Handy said.

Ely greeted Tibor without looking at him. Formal. Correct. Her dislike was careful, almost polite. Father Handy found that worse than open hatred.

Watching Tibor drink coffee was painful.

He locked the cart, rerouted power, extended a metal arm tipped with delicate fingers, adjusted the angle, steadied the tremor. Every simple act had become a technical operation. The coffee sloshed dangerously before finally landing in the cup.

Father Handy lit a cigarette. His body felt wrong today too. Glands misfiring. Hormones out of balance. He wondered if pain was sometimes a warning.

They exchanged German phrases, ritual and comfort. A shared past of books, music, evenings when the light was too bad to paint and all they could do was talk. Culture survived like an ember, fragile and necessary.

Travel no longer connected places. The world had fractured into regions, rumors, static on the radio. Pilgrims left and often never returned.

“You understand,” Father Handy quoted lightly.

“I think I do,” Tibor replied, setting down the cup with exaggerated care.

Silence stretched.

“What is it?” Tibor asked finally.

Father Handy felt the weight of the network he served—commands passed down through belief, obedience disguised as meaning. They weren’t slaves, he told himself. They chose this.

“Why do you work here?” he asked Tibor.

“You pay me,” Tibor said.

“And if I didn’t?”

“I’d starve somewhere else.”

They spoke of power. Of how even gentle authority still compelled. Tibor understood this too well.

Then Father Handy told the story of the ram.

An old animal. Too slow to defend its flock. Forced to watch while a dog killed one of its own.

“If that were you,” Father Handy said quietly, “you would wish you were already dead.”

“Yes,” Tibor said.

Death, in their doctrine, was not an enemy. It was a solution when purpose failed.

“What is my job?” Tibor asked.

Father Handy swallowed.

“To paint Him,” he said. “As He truly is.”

The God of Wrath.

The man blamed for ending the world.

Tibor stared as Father Handy produced a photograph—color, depth, three-dimensional. Pre-war technology. Rare.

The face stared back.

Human. Smiling. Relaxed.

“I can’t paint this,” Tibor said.

Because gods were not supposed to look like men enjoying a meal.

“Then you must find Him,” Father Handy said. “The Church demands it.”

A pilgrimage. Across a broken continent. With failing batteries. With a cow that could die at any time.

“If I fail,” Tibor said, “that’s on you.”

At the stove, Ely muttered, “Fire him.”

“I fire no one,” Father Handy said softly.

He recited an old verse. A fragile bond between two men clinging to meaning.

Tibor drank his coffee.

Outside, the cow shifted and groaned, hungry.

Father Handy thought: if Tibor dies, the mural dies. If the mural dies, belief collapses. No one survives alone anymore.

He wondered if Tibor understood that.

It wouldn’t help.

In this world, nothing did.
